% !TEX program = xelatex
% ============================================================
%  计算机软件著作权说明书 —— 主文件
%
%  本文件包含软件信息配置、文档结构与内容引入；
%  样式配置（宏包、字体、字号、标题格式、参考文献等）
%  集中在 softcopyright.sty。
% ============================================================
\documentclass[12pt]{ctexart}
% 说明：不指定 fontset 选项即为 ctex 自动检测（等效 fontset=auto）：
%   Windows -> 宋体/黑体，macOS -> Songti SC，Linux -> Fandol
% 请勿手写 fontset=windows 等硬编码，否则无法跨平台编译。
\usepackage{softcopyright}

% ------------------------------------------------------------
%  作者与版本信息（请在此处修改）
% ------------------------------------------------------------
% 定义统一的软件名称与版本号
\newcommand{\SoftwareName}{软件名称}
\newcommand{\Version}{V1.0}

% 定义纯作者列表（不含“作者：”，用换行分隔）
\newcommand{\PureAuthors}{%
	作者1 \\
	作者2 \\
	作者3%
}

% 带前缀的完整作者显示（“作者：”不参与替换）
\newcommand{\BasicAuthor}{作者：\\\PureAuthors}

% 修订记录的作者（仅处理纯作者列表，换行转逗号）
\newcommand{\RevisedAuthor}{%
	\StrSubstitute{\PureAuthors}{\\}{, }% 只替换作者间的换行
}

% ------------------------------------------------------------
%  【编写技巧】单独编译某一章以大幅提速：
%  取消下面一行的注释，即只编译“2_operation”一章
%  （配合 \include 使用；\includeonly 参数须与 \include 完全一致）：
%
%   \includeonly{texfile/2_operation}
%
%  注意：正式提交/生成最终 PDF 前，请务必注释掉这一行做全量编译。
% ============================================================

\begin{document}
	% ========== 封面 ==========
	\title{\erhao\textbf{\SoftwareName}\\[0.5em] \sanhao 使用说明}
	% 设置作者和日期
	\author{ \BasicAuthor }
	\date{}
	\maketitle

	% ========== 文件修订记录 ==========
	\newpage
	\phantomsection
	\section*{文件修订记录}
	\addcontentsline{toc}{section}{文件修订记录}  % 手动添加到目录

	% \arraystretch 仅作用于本表格：用 \begingroup ... \endgroup
	% 限定作用域，避免影响文档后续所有表格的行高
	\begingroup
	\renewcommand{\arraystretch}{1.5} % 调整行高为原来的 1.5 倍
	%保持空行不删除，维持合格的表格样式
	\begin{center}
		\begin{tabular}{ | m{1.5cm}<{\centering}|m{4cm}<{\centering} | m{4.64cm}<{\centering} | m{3.5cm}<{\centering} |}
			\hline
			版本号 & 生成日期 & 作者 & 修订内容 \\ % 表头
			\hline
			V1.0 & 2025年6月29日 & \RevisedAuthor  & 初始版本 \\ % 第一行
			\hline
			&  &  &  \\ % 第二行
			\hline
			&  &  &  \\ % 空行
			\hline
			&  &  &  \\ % 空行
			\hline
			&  &  &  \\ % 空行
			\hline
			&  &  &  \\ % 空行
			\hline
			&  &  &  \\ % 空行
			\hline
			&  &  &  \\ % 空行
			\hline
			&  &  &  \\ % 空行
			\hline
			&  &  &  \\ % 空行
			\hline
			&  &  &  \\ % 空行
			\hline
			&  &  &  \\ % 空行
			\hline
			&  &  &  \\ % 空行
			\hline
			&  &  &  \\ % 空行
			\hline
		\end{tabular}
	\end{center}
	\endgroup

	% ========== 目录 ==========
	\newpage
	\phantomsection
	\section*{目录}
	\addcontentsline{toc}{section}{目录}  % 手动添加到目录
	\makeatletter
	\@starttoc{toc}% 直接输出目录内容，避免覆盖 \tableofcontents 原始定义
	\makeatother

	% ========== 图目录 ==========
	\newpage
	\phantomsection
	\section*{图目录}
	\addcontentsline{toc}{section}{图目录}
	\listoffigures

	% ========== 表目录 ==========
	\newpage
	\phantomsection
	\section*{表目录}
	\addcontentsline{toc}{section}{表目录}
	\listoftables

	% ========== 正文 ==========
	% 使用 \include 引入章节（而非 \input）：
	%   1. \include 会在每个章节前后自动分页，符合软著说明书
	%      每章从新页开始的要求（无需再手动写 \newpage）；
	%   2. 配合文件头部的 \includeonly 可单独编译某一章，大幅
	%      提升编写时的编译速度。
	%%%%%%%%%%%%%%
%  第一部分 软件概述
%%%%%%%%%%%%%%
\section{软件概述}

\subsection{项目背景}
% 介绍本软件的开发背景与业务需求。
此处填写项目背景。

\subsection{编写目的}
% 说明本说明书的编制目的与主要作用。
此处填写编写目的。

\subsection{适用对象}
% 说明本软件的目标用户群体。
此处填写适用对象。

\subsection{适用范围}
% 说明本软件支持的运行环境与适用业务范围。
此处填写适用范围。




	%%%%%%%%%%%%%%
%
%  第二部分 操作教程
%
%%%%%%%%%%%%%%
\section{操作教程}
\subsection{操作教程1}

以下为参考文献引用的示例写法，正式提交时请删除：
\upcite{example:article} \upcite{example:book}

	%%%%%%%%%%%%%%
%  第三部分 总结展望
%%%%%%%%%%%%%%
\section{总结展望}

\subsection{总结}
% 对软件的功能实现、运行效果等进行客观总结。
此处填写系统总结。

\subsection{展望}
% 说明软件后续的版本升级与改进设想。
此处填写未来展望。







	% ========== 参考文献 ==========
	% 根据 softcopyright.sty 中启用的方案自动输出：
	%   方案 A（biblatex） -> \printbibliography
	%   方案 B（bibtex）   -> \bibliographystyle + \bibliography
	\clearpage
	\phantomsection
	\addcontentsline{toc}{section}{参考文献}
	\ifcsdef{printbibliography}
		{\printbibliography[title={参考文献}]}   % 方案 A：biblatex + biber
		{\bibliographystyle{gbt7714-numerical}\bibliography{refs}} % 方案 B：bibtex

\end{document}
